\section{Gradient-boosted trees and neural tabular learners}
\label{sec:bg:tabular}

When the inputs are a table of hand-made features, one row per sample, the reference
learner is not a neural network but an ensemble of decision trees built by gradient
boosting.  This section explains the decision tree, the boosting procedure that adds trees
one at a time along the gradient of the loss, the LightGBM implementation used here, and
the two neural learners trained on the same table to tell the learner apart from the
training regime.

\subsection{Decision trees}

A \emph{decision tree} \cite{breiman1984cart} partitions the input space with a sequence
of binary questions on single features, ``is the time since the last move below 17
minutes?'', each question splitting the samples that reach it into two groups, until the
groups, the \emph{leaves}, are small or pure; every leaf then carries one output, the
class frequencies of the training samples in it or, in the trees of a boosted ensemble, a
single number.  A tree is grown greedily: at every node the feature and the threshold are
chosen that most reduce an \emph{impurity} of the two children, the entropy or the Gini
index of their class mix for a classification tree, or the squared error for a regression
tree.  A tree is invariant to any monotone rescaling of the features, needs no
normalisation and no imputation beyond a rule for missing values, and represents sharp
thresholds and interactions between a few features exactly, which an MLP can only
approximate smoothly; a single deep tree, on the other hand, overfits, and a shallow one
underfits.  The depth-three tree of \cref{sec:structure:sep} is a tree of this kind read
in full.

\subsection{Gradient boosting}
\label{sec:bg:boosting}

\emph{Boosting} builds a strong model as a sum of many weak ones, each fitted to what the
previous ones got wrong.  \emph{Gradient boosting} \cite{friedman2001gbm} makes this
precise as gradient descent in the space of functions.  The model after $k$ rounds is an
additive sum of trees,
\begin{equation}
  F_k(\mathbf{x}) = F_{k-1}(\mathbf{x}) + \eta\, h_k(\mathbf{x}), \qquad
  F_0(\mathbf{x}) = \text{a constant (the log-frequencies of the classes)},
  \label{eq:boosting}
\end{equation}
where $h_k$ is a small tree and $\eta$ a \emph{learning rate} (0.05 here) that shrinks
every tree's contribution so that many trees share the work.  For a classifier
$F(\mathbf{x})$ is a vector of three scores, the logits of \cref{eq:softmax}, and one tree
per class is added at every round.  Each round asks which tree would most reduce the loss
$\sum_i \ell(F_{k-1}(\mathbf{x}_i) + h(\mathbf{x}_i), y_i)$, and answers it with a
second-order Taylor expansion of the loss around the current model
\cite{chen2016xgboost}: with
\begin{equation}
  g_i = \frac{\partial \ell}{\partial F}\Big|_{F_{k-1}(\mathbf{x}_i)}, \qquad
  h_i = \frac{\partial^2 \ell}{\partial F^2}\Big|_{F_{k-1}(\mathbf{x}_i)},
\end{equation}
the gradient and the curvature of the loss at every training row (for the cross-entropy
and class $c$, $g_i = \hat{p}_{i,c} - [y_i = c]$ and $h_i = \hat{p}_{i,c}(1 -
\hat{p}_{i,c})$, the same quantities as in \cref{eq:xent-grad}), the loss change produced
by a tree that assigns the value $w_j$ to the rows of its leaf $j$ is approximately
$\sum_j \bigl( G_j w_j + \tfrac{1}{2} (H_j + \lambda) w_j^2 \bigr)$ with $G_j = \sum_{i \in
j} g_i$ and $H_j = \sum_{i \in j} h_i$, which is minimised by
\begin{equation}
  w_j^{\ast} = -\frac{G_j}{H_j + \lambda}, \qquad
  \text{gain of a split} = \frac{1}{2}\Bigl( \frac{G_L^2}{H_L + \lambda} + \frac{G_R^2}{H_R + \lambda} - \frac{(G_L + G_R)^2}{H_L + H_R + \lambda} \Bigr),
  \label{eq:leaf}
\end{equation}
where $\lambda$ is an L2 penalty on the leaf values (1 here).  The tree is grown by
choosing, at every node, the feature and threshold whose split has the largest gain, and
its leaves take the values $w_j^{\ast}$: a leaf's value is a Newton step of the loss,
averaged over the rows in the leaf, so each round moves the scores of every row towards
the direction that reduces its loss, by an amount that depends on how many rows share the
leaf and how sure the model already is.  \Cref{fig:bg:boosting} draws the loop.  Class
weights enter by multiplying $g_i$ and $h_i$ of each row by $\alpha_{y_i}$, and the
offsets of \cref{eq:logit-adjustment} are added to the final scores exactly as for a
network.

\begin{figure}[t]
\centering
\resizebox{\textwidth}{!}{\begin{tikzpicture}[>=Latex, x=1cm, y=1cm,
  box/.style={draw=vizInk, rounded corners=2pt, align=center, font=\footnotesize, inner sep=4pt, minimum height=10mm, text width=36mm, fill=white},
  input/.style={box, dashed, text width=28mm}, op/.style={box, fill=vizB!12}, state/.style={box, fill=vizA!12, text width=30mm}, outbox/.style={box, fill=vizC!15},
  arr/.style={->, draw=vizInk, line width=0.8pt}, lab/.style={font=\scriptsize, text=vizInk, align=center}]
\node[input] (data) at (0,1.2)     {training rows\\$(\mathbf{x}_i, y_i)$, $i = 1, \ldots, N$};
\node[state] (F)    at (0,-1.6)    {current model $F_{k-1}(\mathbf{x})$,\\three scores\\(round 0: class log-frequencies)};
\node[op]    (grad) at (4.8,-0.2)  {for every row, gradient and curvature of the loss\\$g_i = \partial \ell / \partial F$, $\;h_i = \partial^2 \ell / \partial F^2$\\(cross-entropy: $g_i = \hat{p}_i - [y_i = c]$)};
\node[op]    (tree) at (9.6,-0.2)  {fit one small tree per class to the $g_i$:\\split where the gain of \cref{eq:leaf} is largest,\\leaf value $w_j = -G_j / (H_j + \lambda)$};
\node[op]    (add)  at (14.4,-0.2) {add it with a small step\\$F_k = F_{k-1} + \eta\, h_k$, $\ \eta = 0.05$};
\node[outbox] (stop) at (14.4,-3.2) {after each round: score on 12 March;\\stop after 40 rounds without improvement};
\draw[arr] (data) -- (grad); \draw[arr] (F) -- (grad); \draw[arr] (grad) -- (tree); \draw[arr] (tree) -- (add);
\draw[arr] (add) -- (stop);
\draw[arr] (stop.west) -| node[lab, pos=0.25, above] {next round, $k \leftarrow k + 1$} (F.south);
\end{tikzpicture}
}
\caption{One round of gradient boosting.  The gradient and the curvature of the loss are
computed at every training row for the current model; one small tree per class is fitted
to them, with leaf values given by \cref{eq:leaf}; the tree is added with a small step;
the model is scored on the early-stopping day and training stops when the score has not
improved for forty rounds.}
\label{fig:bg:boosting}
\end{figure}

\paragraph{LightGBM.}  LightGBM \cite{ke2017lightgbm} is the implementation used.  Three
of its choices matter here.  It bins every feature into at most 255 \emph{histogram} bins
before training, so that the gain of every candidate split is computed from sums over
bins rather than over rows, which makes a round on seven million rows take seconds.  It
grows trees \emph{leaf-wise}, splitting at every step the leaf with the largest gain
anywhere in the tree rather than completing one level at a time, up to a maximum number of
leaves (63 here) and with a minimum number of rows per leaf (500), which is what limits
how fine a threshold the trees can fit.  And it subsamples the rows and the features of
every tree (80\% of each), so that consecutive trees see slightly different data and the
ensemble does not fit the noise of any one row.  The number of rounds is chosen by
\emph{early stopping} on a day the trees do not train on: the score on 12 March is
computed after every round and training stops when it has not improved for forty rounds
(536 rounds at one hour, hence 1{,}608 trees).  Gradient-boosted trees are the reference
for tabular prediction: the M5 forecasting competition was won almost entirely by LightGBM
on lagged and calendar features \cite{makridakis2022m5}, and systematic comparisons on
tabular benchmarks find trees at or near the top unless a neural model is trained with
care.

\subsection{Neural tabular learners}
\label{sec:bg:tabnets}

That care consists of a few known ingredients \cite{holzmuller2024betterbydefault}.  The
first is the \emph{piecewise-linear encoding} of every numeric feature
\cite{gorishniy2022embeddings}: the range of a feature is cut at $T$ quantiles of the
training data, $b_0 < b_1 < \cdots < b_T$, and the feature value $x$ is replaced by $T$
numbers,
\begin{equation}
  \operatorname{PLE}(x)_t = \min\Bigl(1, \max\Bigl(0, \frac{x - b_{t-1}}{b_t - b_{t-1}}\Bigr)\Bigr), \qquad t = 1, \ldots, T,
  \label{eq:ple}
\end{equation}
which are one for the bins below $x$, zero for the bins above it and a fraction in the bin
that contains it; a linear layer on this encoding can represent any piecewise-linear
function of $x$ with knots at the quantiles, that is, the thresholds that a tree would
find, which is exactly what a plain MLP on the raw feature cannot do.  With $T = 16$ bins
the 59 numeric features become $59 \times 16$ numbers, plus a flag per missing value, the
hour as a sine and a cosine and the category as eight indicators: 674 inputs at one hour.
The other ingredients are the ones of \cref{sec:bg:learning}: shuffled mini-batches of
8{,}192 rows, the one-cycle learning-rate schedule, dropout and early stopping on 12 March.
The MLP so trained has three layers of 512 units.  \emph{TabM} \cite{gorishniy2025tabm}
turns it into a cheap \emph{ensemble}: $k = 16$ members share the same three layers, but
each member multiplies the inputs of every layer by its own learned vector (initialised at
$\pm 1$) and has its own output head, so that the sixteen members are different functions
at the cost of one; the prediction is the mean of their sixteen softmax outputs.  Averaging
several differently initialised classifiers is the most reliable known improvement of a
neural classifier, and TabM obtains most of it with a tenth of the parameters of sixteen
separate networks; it is among the best neural tabular learners on recent benchmarks
\cite{jiang2026omnitabbench}.  Both learners are used in \cref{ch:stateless} for one
purpose: to tell whether a gap between the trees and the stateful model is the learner or
the training regime.
